% =============================================================================
% Visual Atlas: Mindmap + Cheatsheet macros (2 landscape A4 pages per leaf)
%   Page 1  = MINDMAP    : root pictogram + 6 colour-coded branches + hook + trap
%   Page 2  = CHEATSHEET : 3-column cards (formulas, recipe, shapes, trade-offs)
% Usage (see any *_cheatsheet.tex):
%   \csbegin{<domain 1-7>}{<Parent topic>}{<Leaf title>}{<tagline>}
%     \begin{tikzpicture} \csroot{..}{..} \csbranch{1..6}{A..G}{..}{..} \csnote{7|8}{..} \end{tikzpicture}
%     \csfoot{..}
%   \csnextpage
%     \begin{cscolumns} \begin{cscard}{A}{Title} .. \end{cscard} \columnbreak .. \end{cscolumns}
%   \csend
% =============================================================================

\usepackage{multicol}
\usepackage{pdflscape}

% --- Branch palette (vivid, high contrast, colour-blind-friendly mix) ---------
\definecolor{csA}{HTML}{E11D48} % crimson  : WHAT / definition
\definecolor{csB}{HTML}{EA580C} % orange   : HOW / mechanism
\definecolor{csC}{HTML}{CA8A04} % amber    : KEY RESULT / theorem
\definecolor{csD}{HTML}{16A34A} % green    : CONNECTIONS
\definecolor{csE}{HTML}{0891B2} % cyan     : COST / complexity
\definecolor{csF}{HTML}{4F46E5} % indigo   : USE-CASES
\definecolor{csG}{HTML}{9333EA} % purple   : hook / trap

% --- Domain colours (title bar) ----------------------------------------------
\definecolor{csDom1}{HTML}{16324F}
\definecolor{csDom2}{HTML}{0E7490}
\definecolor{csDom3}{HTML}{3730A3}
\definecolor{csDom4}{HTML}{7E22CE}
\definecolor{csDom5}{HTML}{1D4ED8}
\definecolor{csDom6}{HTML}{15803D}
\definecolor{csDom7}{HTML}{C2410C}
\colorlet{csDom}{csDom1}

% --- Lists inside boxes -------------------------------------------------------
\newlist{csbul}{itemize}{1}
\setlist[csbul]{label=\textcolor{csDom}{\raisebox{0.1ex}{\scriptsize$\blacktriangleright$}},
  leftmargin=1.0em, itemsep=0.7pt, topsep=0.5pt, parsep=0pt, partopsep=0pt}

% --- Page title bar -----------------------------------------------------------
% #1 kicker (MINDMAP / CHEATSHEET)  #2 parent topic  #3 title  #4 tagline
\newcommand{\cstitlebar}[4]{%
  \noindent\begin{tikzpicture}
    \fill[csDom, rounded corners=4pt] (0,0) rectangle (\textwidth,1.25cm);
    \fill[white, opacity=0.16, rounded corners=4pt] (0,0) rectangle (3.0cm,1.25cm);
    \node[white, font=\sffamily\bfseries\small, anchor=center] at (1.5cm,0.625cm) {#1};
    \node[white, font=\sffamily\bfseries\Large, anchor=west, text width=\dimexpr\textwidth-3.6cm\relax]
      at (3.3cm,0.80cm) {#3};
    \node[white, font=\sffamily\scriptsize, anchor=west, text width=\dimexpr\textwidth-3.6cm\relax]
      at (3.3cm,0.30cm) {\textbf{#2}\;\textbar\; #4};
  \end{tikzpicture}\par\nointerlineskip\vspace{0.15cm}%
}

% --- Page lifecycle -----------------------------------------------------------
% Landscape via pdflscape: the page stays A4-portrait in the PDF with /Rotate 90,
% so XeLaTeX (xdvipdfmx) can mix portrait chapters and landscape atlas pages.
% #1 domain index  #2 parent topic  #3 leaf title  #4 tagline
\newcommand{\csbegin}[4]{%
  \colorlet{csDom}{csDom#1}%
  \gdef\csParent{#2}\gdef\csTitle{#3}\gdef\csTagline{#4}%
  \newgeometry{margin=0.7cm,headheight=0pt,headsep=0pt,footskip=0pt}%
  \landscape
  \setlength{\textwidth}{\dimexpr\paperheight-1.4cm\relax}%
  \setlength{\textheight}{\dimexpr\paperwidth-1.4cm\relax}%
  \setlength{\hsize}{\textwidth}\setlength{\linewidth}{\textwidth}%
  \setlength{\oddsidemargin}{-1in}\setlength{\evensidemargin}{-1in}%
  \setlength{\topmargin}{-1in}\setlength{\headheight}{0pt}\setlength{\headsep}{0pt}%
  \addtolength{\oddsidemargin}{0.7cm}\addtolength{\evensidemargin}{0.7cm}%
  \addtolength{\topmargin}{0.7cm}%
  \thispagestyle{empty}\phantomsection
  \addcontentsline{toc}{chapter}{Atlas: #3}%
  \cstitlebar{MINDMAP}{#2}{#3}{#4}%
  \centering
}
\newcommand{\csnextpage}{%
  \par\newpage\thispagestyle{empty}\raggedright
  \cstitlebar{CHEATSHEET}{\csParent}{\csTitle}{\csTagline}%
}
\newcommand{\csend}{\par\endlandscape\restoregeometry}
\newcommand{\csfoot}[1]{\par\vspace{-0.1cm}{\centering\sffamily\tiny\color{monographSlate} #1\par}}

% --- Mindmap geometry ---------------------------------------------------------
\tikzset{
  csroot/.style={circle, draw=csDom, fill=white, line width=3pt, align=center,
    inner sep=4pt, text width=5.0cm, font=\sffamily\bfseries\normalsize, text=csDom},
}
\pgfdeclarelayer{csbg}
\pgfsetlayers{csbg,main}

% slot -> position & anchors. Slots 1-3 left, 4-6 right, 7 top-centre, 8 bottom-centre.
\newcommand{\csslot}[1]{%
  \ifcase#1\or
    \def\csx{-9.75}\def\csy{5.86}\def\csedge{-5.48}\def\csout{152}\def\csin{0}\or
    \def\csx{-9.75}\def\csy{0}\def\csedge{-5.48}\def\csout{180}\def\csin{0}\or
    \def\csx{-9.75}\def\csy{-5.86}\def\csedge{-5.48}\def\csout{208}\def\csin{0}\or
    \def\csx{9.75}\def\csy{5.86}\def\csedge{5.48}\def\csout{28}\def\csin{180}\or
    \def\csx{9.75}\def\csy{0}\def\csedge{5.48}\def\csout{0}\def\csin{180}\or
    \def\csx{9.75}\def\csy{-5.86}\def\csedge{5.48}\def\csout{-28}\def\csin{180}\or
    \def\csx{0}\def\csy{7.18}\def\csedge{0}\def\csout{90}\def\csin{90}\or
    \def\csx{0}\def\csy{-7.18}\def\csedge{0}\def\csout{-90}\def\csin{-90}\fi
}

\newtcolorbox{csboxenv}[3][]{enhanced, colback=#2!6, colframe=#2, coltitle=white,
  colbacktitle=#2, fonttitle=\sffamily\bfseries\footnotesize, title={#3},
  width=8.55cm, height=5.60cm, boxrule=1.1pt, arc=2.5mm, top=1.5pt, bottom=1.5pt,
  left=4pt, right=4pt, toptitle=1pt, bottomtitle=1pt,
  halign=left, valign=top, fontupper=\sffamily\fontsize{7.5pt}{8.7pt}\selectfont,
  before upper={\raggedright\setlength{\tabcolsep}{2.2pt}\setlength{\parskip}{0.4pt}}, #1}

\newtcolorbox{csnoteenv}[3][]{enhanced, colback=#2!10, colframe=#2, coltitle=white,
  colbacktitle=#2, fonttitle=\sffamily\bfseries\scriptsize, title={#3},
  width=9.3cm, height=2.95cm, boxrule=1.1pt, arc=2.5mm, top=1.2pt, bottom=1.2pt,
  left=4pt, right=4pt, toptitle=0.8pt, bottomtitle=0.8pt,
  halign=left, valign=top, fontupper=\sffamily\fontsize{7.5pt}{8.7pt}\selectfont,
  before upper={\raggedright\setlength{\parskip}{0.4pt}}, #1}

% Root: #1 title text  #2 pictogram (TikZ picture, ~ 3.0cm tall)
\newcommand{\csroot}[2]{%
  \node[csroot] (root) at (0,0) {#1\\[2pt]{\normalfont\footnotesize\tikzset{shape=rectangle, line width=0.4pt, text width=, text=black, inner sep=1.5pt}#2}};%
}

% Branch: #1 slot 1..6  #2 colour letter A..G  #3 header  #4 body
\newcommand{\csbranch}[4]{%
  \csslot{#1}%
  \begin{pgfonlayer}{csbg}
    \draw[cs#2, line width=2.8pt, -{Circle[length=5pt]}, line cap=round]
      (0,0) to[out=\csout,in=\csin] (\csedge,\csy);
  \end{pgfonlayer}%
  \node[inner sep=0pt] at (\csx,\csy) {\begin{csboxenv}{cs#2}{#3}#4\end{csboxenv}};%
}

% Note banner (slots 7, 8): #1 slot  #2 colour  #3 header  #4 body
\newcommand{\csnote}[4]{%
  \csslot{#1}%
  \node[inner sep=0pt] at (\csx,\csy) {\begin{csnoteenv}{cs#2}{#3}#4\end{csnoteenv}};%
}

% --- Cheatsheet page cards ----------------------------------------------------
\newtcolorbox{cscard}[3][]{enhanced, colback=#2!5, colframe=#2, coltitle=white,
  colbacktitle=#2, fonttitle=\sffamily\bfseries\small, title={#3},
  boxrule=1pt, arc=2mm, top=1.5pt, bottom=1.5pt, left=4pt, right=4pt,
  toptitle=1pt, bottomtitle=1pt, fontupper=\sffamily\small,
  before upper={\raggedright\setlength{\tabcolsep}{3pt}\setlength{\parskip}{1.5pt}},
  before skip=2.5pt, after skip=2.5pt, nobeforeafter, #1}
\newenvironment{cscolumns}{%
  \setlength{\columnsep}{7pt}\setlength{\columnseprule}{0pt}%
  \setlength{\parskip}{0pt}\begin{multicols}{3}\raggedcolumns}{\end{multicols}}

% Tiny helpers for compact math in cards
\newcommand{\csfm}[1]{\par\vspace{0.5pt}\noindent\hfil$\displaystyle #1$\hfil\par\vspace{0.5pt}}
\newcommand{\cshl}[1]{\textbf{\textcolor{csDom}{#1}}}
